\documentclass[11pt]{article}


\setlength{\parindent}{0pt}
\usepackage{xltxtra}
\usepackage{hyperref}
\hypersetup{hidelinks}
\usepackage{url}
\urlstyle{tt}
\usepackage{xcolor}
\definecolor{CVBlue}{RGB}{23,110,191}
\definecolor{SoftGray}{RGB}{105,105,105}
\usepackage{calc}
\usepackage{graphicx}
\usepackage{fontspec}
\usepackage{xeCJK}
\usepackage{enumitem}
\CJKsetecglue{}


\setmainfont[
    Path = fonts/Main/,
    Extension = .otf,
    BoldFont = texgyretermes-bold.otf
]{texgyretermes-regular.otf}


\setCJKmainfont[
    Path = fonts/hansans/,
    Extension = .ttf,
    BoldFont = NotoSansSC-Bold.ttf
]{NotoSansSC-Regular.ttf}


\usepackage{fontawesome}
\usepackage[
    a4paper,
    left=1.3cm,
    right=1.3cm,
    top=1.0cm,
    bottom=0.8cm,
    nohead
]{geometry}


\renewcommand{\baselinestretch}{1.15}


\usepackage{titlesec}
\setlist{noitemsep}
\setlist[itemize]{topsep=0.12em, itemsep=0.06em, leftmargin=1.1em}


\newlength{\interProjectSpacing}
\setlength{\interProjectSpacing}{0.34em}
\newcommand{\projectsep}{\vspace{\interProjectSpacing}}


\newlength{\intraProjectTitleSep}
\setlength{\intraProjectTitleSep}{0.16em}
\newcommand{\titlebreak}{\\[\intraProjectTitleSep]}


\newlength{\intraProjectListTopSep}
\setlength{\intraProjectListTopSep}{0.05em}


\titleformat{\section}
{\large\bfseries\raggedright}
{}{0em}
{}
[{\color{CVBlue}\titlerule}]
\titlespacing*{\section}{0cm}{*0.85}{*0.45}


\newcommand{\infoRow}[2]{%
    \textbf{#1} & #2 \\[0.22em]
}


\begin{document}
\pagenumbering{gobble}


\noindent
\begin{minipage}[c]{0.27\textwidth}
    \includegraphics[height=4.4cm]{zju.png}
\end{minipage}
\hfill
\begin{minipage}[c]{0.70\textwidth}
    \raggedright
    {\fontsize{19}{22}\selectfont\bfseries 胡光望}\\[0.18em]
    {\fontsize{11.5}{14}\selectfont 浙江大学能源工程学院车辆工程}\\[0.30em]
    {\fontsize{10}{12.5}\selectfont 2004.10 \quad 重庆彭水 \quad 共青团员}\\[0.08em]
    {\fontsize{10}{12.5}\selectfont 19557148649 \quad dhu977668@gmail.com}
\end{minipage}


\par\vspace{0.35em}
{\color{SoftGray}\rule{\textwidth}{0.35pt}}
\vspace{0.15em}


\section{\makebox[\widthof{\faInfo}][c]{\color{CVBlue}\faInfo}\ 个人概述}
\begin{itemize}[nosep]
    \item 浙江大学能源工程学院车辆工程本科在读，\textbf{主修绩点 4.21}。课程基础覆盖车辆动力、工程热力学、传热学、电工电子和自动控制。
    \item 具备较强的自主学习与学科交叉意识，能够自主学习并使用 \textbf{Python、C、CARLA、Simulink、LabVIEW} 等工具，将机械、能源、控制和数据采集问题串联起来完成方案设计与验证。
\end{itemize}


\section{\makebox[\widthof{\faGraduationCap}][c]{\color{CVBlue}\faGraduationCap}\ 教育背景}
\textbf{浙江大学} \hfill 2023.09 -- 至今\\[0.22em]
能源工程学院 \quad 车辆工程 \quad 本科\\[0.32em]
\renewcommand{\arraystretch}{1.08}
\begin{tabular}{@{}p{2.5cm}p{14.8cm}@{}}
    \textbf{车辆与动力} & 汽车构造、汽车理论、汽车设计、汽车动力系统原理、新能源汽车技术 \\
    \textbf{热能与流体} & 工程热力学（甲）、传热学（甲） \\
    \textbf{控制与电子} & 自动控制理论、汽车电子与控制、电工电子学及实验 \\
    \textbf{机械与实践} & 设计与制造 I/II/III、电子工程训练 \\
\end{tabular}


\section{\makebox[\widthof{\faUsers}][c]{\color{CVBlue}\faUsers}\ 项目经历}


\textbf{基于 CARLA 的雨夜行人横穿避让场景仿真} \hfill 课程项目 \titlebreak
\begin{itemize}[nosep, topsep=\intraProjectListTopSep]
    \item 构建夜间强降雨、前车急刹、行人横穿叠加场景，配置湿滑路面、雾灯和固定地图点位，提升实验复现性。
    \item \textbf{控制策略：} 基于语义 LiDAR 判断风险，设计“紧急左变道--回正--制动”状态机，对比分析 \textbf{PID 与模糊控制} 在横纵向避让优化中的适用性。
\end{itemize}


\projectsep


\textbf{基于水力测功器的发动机台架测控系统设计} \hfill 课程项目 \titlebreak
\begin{itemize}[nosep, topsep=\intraProjectListTopSep]
    \item 设计“发动机--水力测功器--测控柜--上位机”总体方案，规划多类传感器、\textbf{STM32} 采集和 CAN/RS-485 通信。
    \item 提出前馈加位置式 \textbf{PID} 的定扭矩控制方案，设计积分限幅、微分滤波、速率限制及 Simulink 建模思路。
\end{itemize}


\projectsep


\textbf{基于 Simulink 的进气道湿壁油膜建模与前馈补偿} \hfill 课程项目 \titlebreak
\begin{itemize}[nosep, topsep=\intraProjectListTopSep]
    \item 建立进气道湿壁油膜 \textbf{Simulink} 模型，分析附壁系数 \(X\) 和时间常数 \(\tau\) 对燃油动态传输的影响。
    \item 设计前馈补偿器使入缸燃油贴合目标；提出基于压力、转速和壁温的查表自适应补偿方向。
\end{itemize}


\projectsep


\textbf{基于 NI cDAQ 与 LabVIEW 的共轨压力/转速同步采集} \hfill 自主实践 \titlebreak
\begin{itemize}[nosep, topsep=\intraProjectListTopSep]
    \item 基于 NI cDAQ-9174、9205 和 9401 搭建压力/转速双通道 LabVIEW 采集程序，完成任务管理与动态显示。
    \item 使用移位寄存器与噪声模型模拟转速惯性，解决仿真数字通道动态响应不足的问题。
\end{itemize}


\section{\makebox[\widthof{\faCogs}][c]{\color{CVBlue}\faCogs}\ 专业技能}
\renewcommand{\arraystretch}{1.10}
\begin{tabular}{@{}p{2.6cm}p{14.7cm}@{}}
    \textbf{编程语言} & Python、C \\
    \textbf{建模与仿真} & MATLAB/Simulink、CARLA、MWORKS \\
    \textbf{测控与电路} & NI LabVIEW、NI Multisim \\
    \textbf{设计与办公} & CAD、WPS \\
    \textbf{AI 工具} & Codex、DeepSeek Harness \\
    \textbf{语言能力} & CET-6 474 \\
\end{tabular}


\end{document}
